\section{Instrumentation and benchmark methodology}
\label{sec:method}

\subsection{Inside the node}

The node keeps lock-free counters (relaxed atomics), including:
\begin{itemize}
  \item frames in and out, and accepts;
  \item send failures and backpressure drops;
  \item reconnects and evictions;
  \item queued sends and flushes;
  \item high-water marks of the render, store and sender queues.
\end{itemize}
With \code{-{}-metrics}, a metrics thread writes one JSONL snapshot per second. Each
snapshot adds \code{getrusage(2)} figures (user and system CPU time, voluntary and
involuntary context switches, maximum RSS) and thread and fd counts from
\code{/proc/self}. Every node thread is named, so a hang dump identifies each thread.

\paragraph{CPU usage.} \code{getrusage(RUSAGE\_SELF)} returns the user time
$u$ and system time $s$ consumed so far by \emph{all} threads of the process. Given
two snapshots taken at wall-clock times $t_a$ and $t_b$, the node's CPU usage over
that interval is
\begin{equation}
\mathrm{CPU} = 100\% \times \frac{(u_b + s_b) - (u_a + s_a)}{t_b - t_a}.
\label{eq:cpu}
\end{equation}
The numerator is CPU time and the denominator is elapsed time, so the ratio is the
average number of cores the process kept busy: 100\% is one core fully used. The node
is pinned to two CPUs, so its ceiling is 200\%. This is the usual way to compute a
process's CPU usage on Linux, whether the times come from \code{getrusage} or from the
\code{utime} and \code{stime} fields of \code{/proc/<pid>/stat}~\cite{poll,so-cpu}.
$t_a$ and $t_b$ are the snapshots that bound the measurement window.

\paragraph{Context switches.} The same \code{getrusage(RUSAGE\_SELF)} call returns
two counters, summed by the kernel over all threads of the process:
\begin{itemize}
  \item \code{ru\_nvcsw}, \emph{voluntary} context switches: a thread gave up the CPU
        itself because it had to wait, for example blocking in \code{recv},
        \code{poll}, \code{epoll\_wait}, on a mutex or on a queue's condition
        variable;
  \item \code{ru\_nivcsw}, \emph{involuntary} context switches: the scheduler took
        the CPU away from a thread that was still runnable, because its time slice
        ran out or another thread had to run.
\end{itemize}
Both counters only grow, so they are turned into rates in the same way as
Equation~\ref{eq:cpu}:
\begin{equation}
\text{voluntary rate} = \frac{\mathit{nvcsw}_b - \mathit{nvcsw}_a}{t_b - t_a},
\qquad
\text{involuntary rate} = \frac{\mathit{nivcsw}_b - \mathit{nivcsw}_a}{t_b - t_a}.
\end{equation}
A high voluntary rate means threads often sleep and wake, for example once per
message. At 1,024 peers the thread modes switch voluntarily 89--107 thousand times
per second, against 6--16 thousand for the modes with one I/O thread. A high
involuntary rate means there were more runnable threads than CPUs.

\subsection{The load generator}
\label{sec:gen}

\code{stressgen} is a single-threaded \code{epoll} client. It simulates many peers that
speak the same protocol as the node, in either link shape (\code{-{}-duplex}).
\begin{itemize}
  \item \textbf{RTT:} each message embeds its send time, and RTT is measured when the
        node's echo arrives, on the same clock. Times are recorded in a log-bucketed
        histogram with 32 sub-buckets per power of two, so a quantile is
        over-estimated by less than 3.125\%.
  \item \textbf{What counts:} only messages sent inside the measurement window, with
        a short grace period for late echoes.
  \item \textbf{Load:} open loop (a fixed \code{-{}-rate} per peer) or closed loop (a
        \code{-{}-window} of outstanding messages per peer).
  \item \textbf{Peer types:} besides active peers, it can run idle peers and slow
        peers (receive buffer 4\,KiB, never read).
  \item \textbf{Churn:} at a fixed rate (\code{-{}-churn}, reconnects per second), the
        generator closes a randomly chosen \emph{ready} peer and redials it at once.
        The schedule is open loop: it does not wait for earlier redials to finish. If
        no peer is ready when a churn event is due, the event is skipped.
  \item \textbf{Handshakes:} it records handshake latency from TCP connect to Ready.
\end{itemize}

\subsection{The harness}

\code{run\_matrix.py} runs the experiment as a matrix of \emph{cells}. A cell is one
scenario, one swept value, one mode and one repetition. For each cell it:
\begin{enumerate}
  \item starts the node in headless echo mode with the chosen \code{-{}-io} mode and
        \code{-{}-metrics}, pinned with \code{taskset} to CPUs 0--1, and waits for the
        node's \code{ready} event on its JSON output;
  \item starts \code{stressgen}, pinned to CPUs 4--7, which dials every simulated peer
        (in batches, so the accept backlog does not overflow) and waits until all of
        them have completed the handshake or a ramp deadline has passed;
  \item lets the generator run a 5\,s warm-up and then a 20\,s measurement window,
        after which it writes its result file;
  \item sends \code{/quit} to the node on its standard input and waits up to 30\,s for
        the reply \code{bye} and for the process to exit. A node that does not exit
        has the state of every thread recorded from \code{/proc} and is then killed;
  \item writes a metadata file with the cell's parameters, the build stamp (commit,
        uncommitted changes, binary hash, build type), the node's exit status and the
        result of the validity checks below.
\end{enumerate}
The CPUs 0--1 and 4--7 are separate physical cores of an Intel Core Ultra~7 155H,
under WSL2 (Linux 6.6). The order of the modes is shuffled for every (scenario, value,
repetition), so slow drift on the machine is spread across modes rather than biasing
one. The script refuses non-Release builds, and when an interrupted run is resumed it
refuses to mix cells from different builds.

Table~\ref{tab:scenarios} lists the scenarios.

\begin{table}[h]
\centering\small
\begin{tabularx}{\linewidth}{@{}llXl@{}}
\toprule
Scenario & Swept & Other load & Network \\
\midrule
S1 fan-in & active peers 1, 8, 64, 256, 1024 & 100 msg/s per peer, 64\,B & loopback \\
S2 idle-heavy & idle peers 0, 100, 300, 1000 & 8 active $\times$ 500 msg/s & loopback \\
S3 slow consumers & slow peers 0, 1, 4 & 16 active $\times$ 200 msg/s & loopback \\
S4 churn & reconnects/s 0, 10, 50 & 64 peers $\times$ 50 msg/s & loopback \\
S5 throughput & size 32\,B, 4\,KiB, 64\,KiB & 8 peers, closed loop, window 16 & loopback \\
S6a WAN churn & RTT 0, 4, 20, 30, 40, 50\,ms & 32 $\times$ 100 msg/s, 256\,B, 16 reconnects/s & veth + netem \\
S6c WAN churn, half & RTT 50\,ms & as S6a, 8 reconnects/s & veth + netem \\
S6b WAN bottleneck & 1000, 200, 50\,Mbit/s & 32 peers, closed loop, window 8, 16\,KiB & veth, 4\,ms RTT \\
S7 capped fan-in & \code{-{}-max-peers} 64, 128, 256, 1040 & 1024 peers $\times$ 100 msg/s & loopback \\
\bottomrule
\end{tabularx}
\caption{Scenarios: 32 swept values $\times$ 9 variants $\times$ 3 repetitions = 864 cells.}
\label{tab:scenarios}
\end{table}

\paragraph{WAN lab.} The S6 cells run inside an unprivileged \code{unshare}
namespace. The node and the generator are joined by a veth pair with MTU 1500 and
offloads disabled. netem~\cite{netem} delay is split evenly between the two
directions, so the configured RTT is their sum. The S6b rate limit applies only to
the node's egress.

\subsection{Analysis and validity}

\code{analyze.py} aligns the node's metrics window with the generator's monotonic
window bounds. From that window it computes CPU usage (Equation~\ref{eq:cpu}) and
context-switch rates per cell, and it reports the median and inter-quartile range over the valid
repetitions.

A cell counts as \emph{degraded}, and is left out of the medians, if any of these
checks fails:
\begin{itemize}
  \item not every peer became ready (at least 90\% with churn; the cap rule in S7);
  \item an open-loop cell offered less than 98\% of its nominal load;
  \item there were connect failures (outside S7), bad frames or skipped sends;
  \item the node stalled: its \code{frames\_in} stopped growing within 3\,s of the
        window's end.
\end{itemize}
Overload is a result, not a failure. A cell that delivered a fraction of its load with
huge latency is reported as the point where that mode breaks.

\subsection{The measured run}

All results come from one run of the full matrix: 864 cells from one clean Release
build (\code{-O3}). Every cell completed, and every node exited with status 0 and no
hang. 840 cells pass every validity check. The 24 degraded cells are all simplex S6a
cells at RTT 40 and 50\,ms (Section~\ref{sec:res-wan}). The build and the location of
the raw data are listed in Appendix~\ref{app:repro}.

\subsection{Limitations}

\begin{itemize}
  \item The node and the generator share one WSL2 host, and loopback RTT includes the
        generator's own processing time.
  \item CPU and context switches are measured for the node only, without its softirq
        share. Threads and fds are sampled once a second.
  \item The node has two CPUs. That is where the thread-per-connection ceiling falls;
        with more cores it would move up.
  \item There is no counter for outbound bytes.
  \item Handshake statistics only include peers that completed a handshake, so they
        carry survivorship bias.
  \item The generator's slow peers never made the duplex shapes queue writes
        (Section~\ref{sec:res-slow}).
\end{itemize}
